%!TEX root = ../main.tex
\chapter{Diseño del sistema de control}

\section{Modelo matemático del sistema}

A lo largo de este capítulo se utilizan dos notaciones para la derivada
temporal: la notación de punto, $\dot{x}$, se emplea al escribir
ecuaciones de estado y sistemas dinámicos en forma compacta (por ejemplo,
$\dot{x} = f(x)$); la notación de fracción, $\frac{dx}{dt}$, se emplea
al despejar o manipular explícitamente la derivada de una variable individual
dentro de un desarrollo algebraico.


\subsection{Subsistema mecánico}

Este subsistema está compuesto por el rotor, el tren motriz y el generador 
eléctrico. En turbinas eólicas convencionales, el tren motriz incluye un 
eje de baja velocidad conectado al rotor, una caja de engranes 
(\emph{gearbox}) y un eje de alta velocidad que alimenta al generador 
\cite{soltani2013}. La caja de engranes convierte la velocidad rotacional 
baja del rotor a una velocidad mayor requerida por el generador, con una 
relación de engranes fija $N_g$. Sin embargo, en configuraciones que 
utilizan un generador síncrono de imanes permanentes (PMSG) de 
accionamiento directo (\emph{direct-drive}), como la considerada en esta 
tesis, el generador opera a la misma velocidad que el rotor, lo que 
elimina la necesidad de la caja de engranes ($N_g = 1$). Esta 
configuración reduce las pérdidas mecánicas, los costos de mantenimiento y 
el peso del conjunto, a cambio de requerir un generador de mayor diámetro 
y mayor número de polos \cite{Burton2011}.
\\

La dinámica del sistema mecánico, válida para un modelo de una masa 
\cite{soltani2013}, está dada por:

\begin{equation}
    J \dot{\omega}_m = -f_r \omega_m + T_m - T_e
    \label{eq:torque_balance_simplified}
\end{equation}

donde $\omega_m$ representa la velocidad angular del rotor, $f_r$ es el 
coeficiente de fricción viscosa, $T_m$ es el par mecánico aplicado por el 
viento y $T_e$ el par eléctrico del generador.
\\

Es importante aclarar que en la práctica la potencia real que se puede obtener del viento es limitada y, por 
tanto, también lo es el torque aerodinámico/mecánico. En particular, la potencia extraída del viento está dada por
\footnote{Esta es una ecuación fundamental en aerodinámica de turbinas eólicas, ampliamente documentada en 
la literatura. Para una derivación detallada véase Burton et al. \cite{Burton2011}. Para aplicaciones prácticas 
y consideraciones de diseño véase Hau \cite{Hau2013}.}:
\begin{equation}
    P = \frac{1}{2}\rho \pi r^2 v_w^3 C_p(\lambda,\beta)
    \label{eq:aerodynamic_power}
\end{equation}

donde $\rho$ es la densidad del aire, $r$ es el radio de giro de las aspas, 
$v_w$ la velocidad del viento y $C_p(\lambda,\beta)$ es el coeficiente de 
potencia en términos de la relación de velocidad punta ($\lambda$) y el 
ángulo de ataque de las aspas ($\beta$).
\\

La potencia se relaciona directamente con el torque aerodinámico mediante la ecuación:
\begin{equation}
    T_r = \frac{P}{\omega_r}.
    \label{eq:aerodynamic_torque}
\end{equation}

En (\ref{eq:torque_balance_simplified}) utilizamos una versión simplificada del balance de torques. 
En esta ecuación se define cada par como sigue:

\begin{equation}
    T_m = r\frac{\kappa_1}{v_w} \frac{C_p(\lambda,\beta)}{\lambda}
    \label{eq:mechanical_torque}
\end{equation}

donde $\kappa_1$ es un parámetro que agrupa la potencia aerodinámica disponible y se define como: 

\begin{equation}
    \kappa_1 = \frac{1}{2}\rho\pi r^2 v_w^3.
\end{equation}

El resto de la ecuación surge a partir de la manipulación algebráica y posterior simplificación de:

\begin{equation}
    \lambda = \frac{r\omega_m}{v_w}.
    \label{eq:lambda}
\end{equation}

Por su parte, el par eléctrico se define como:

\begin{equation}
    T_e = \frac{3}{2}p\phi i_q.
    \label{eq:electrical_torque}
\end{equation}

Estas ecuaciones introducen el flujo eléctrico en el generador ($\phi$) y la corriente en el eje de 
cuadratura del generador ($i_q$). Estas últimas dos variables serán tratadas
a profundidad en la sección del subsistema eléctrico.


\subsection{Coeficiente de potencia}\label{sec:coef_potencia}

Durante su operación, el rotor de la turbina eólica absorbe energía de la corriente de aire, 
lo cual promueve su movimiento. El poder de absorción y las condiciones de operación del aerogenerador dependen, por tanto, 
del área efectiva de barrido ($A_R := \pi r^2$), la velocidad del viento y los cambios que ocurren en estas 
cantidades en el campo de flujo del rotor. Durante el proceso de extracción de la potencia el flujo de aire experimenta
desaceleración axial, una desviación tangencial al momento de chocar con el área del rotor y una 
expansión del área de la sección transversal una vez que ha pasado el rotor (efecto estela) \cite{heier2014}.
\\

La potencia teórica extraída por la turbina se puede expresar como:

\begin{align}
    P_W = A_R\frac{\rho}{2}v^3
    \label{eq:theoretical_power}
\end{align}

y de acuerdo con Betz \cite{betz1926}, la potencia máxima teórica extraída sería:

\begin{align}
    P_{W_{\text{max}}} = \frac{16}{27}A_R\frac{\rho}{2}v^3 
    \label{eq:max_theoretical_power}
\end{align}

El cociente entre la potencia teórica extraída y la potencia en la masa de aire en movimiento, la cual se expresa como:

\begin{align}
    P_0 = A_R\frac{\rho}{2}v_1^3
\end{align}

define el coeficiente de potencia (también llamado coeficiente de rendimiento):

\begin{align}
    C_p = \frac{P_W}{P_0} = \frac{16}{27}\frac{v^3}{v_1^3}
\end{align} 

Se puede observar que el coeficiente de potencia consiste en el cociente entre el cubo de la velocidad $v$ 
(el cual es una función de las velocidades $v_1$, $v_2$ y $v_3$, siendo estas las velocidades antes, en 
y después del rotor, respectivamente), y el cubo de la velocidad del viento antes del rotor ($v_1$). En 
la práctica, la manera más común de obtener este valor es utilizando métodos 
computacionales y artificios matemáticos más complejos (se puede encontrar un desarrollo detallado del 
método \emph{blade element} en \cite{heier2014}).
\\

Un punto importante a notar es que, de acuerdo con \cite{Burton2011}, en turbinas modernas de tres palas 
el valor máximo de $C_p$ es de solo 0.47, lo cual se aleja considerablemente del límite teórico de Betz, 
y se logra con un valor de la relación de velocidad punta de $\lambda^\star \approx 7$. Este valor de 
$\lambda^\star$ corresponde a un valor típico reportado en la literatura para turbinas comerciales de 
tres palas, y no debe confundirse con el valor de $\lambda^\star$ calculado más adelante en esta sección 
para el modelo específico de $C_p(\lambda,\beta)$ utilizado en esta tesis.
\\

Con el objetivo de obtener un valor del coeficiente de potencia ajustado a los parámetros utilizados durante 
el modelado de la turbina se utilizará una aproximación mediante una función analítica de $C_p$.
En \cite{man1981} se desarrolla una expresión para el coeficiente de potencia como la siguiente función de 
la relación de velocidad punta y el ángulo de inclinación de las palas:

\begin{align}
    C_p = c_1 \left( c_2 - c_3\beta - c_4\beta^x - c_5 \right) e^{-c_6(\lambda,\beta)}
    \label{eq:cp_analytic}
\end{align}

Los valores de las constantes cambian según el método y modelo utilizado. A partir del procedimiento en 
\cite{amlang1992}, los valores de estas constantes se definen como:
\begin{align}
    \begin{array}{lll}
        c_1 = 0.5, &c_2 = \frac{116}{\lambda_i}, &c_3 = 0.4,
        \\
        c_4 = 0, &c_5 = 5, &c_6 = \frac{21}{\lambda_i}
    \end{array}
    \label{eq:cp_constants}
\end{align}
donde $\lambda_i$ está dado por
\begin{align}
    \frac{1}{\lambda_i} = \frac{1}{\lambda + 0.08\beta} - \frac{0.035}{\beta^3 + 1}.
    \label{eq:lambda_i}
\end{align}

Nótese que, si bien $c_6$ se expresa únicamente en términos de $\lambda_i$, esta última depende tanto de 
$\lambda$ como de $\beta$, por lo que la notación $c_6(\lambda,\beta)$ es consistente. Por otro lado, 
\cite{amlang1992} no reporta un valor explícito para el exponente $x$ en este conjunto de constantes; sin 
embargo, dado que $c_4 = 0$, el término $c_4\beta^x$ se anula independientemente del valor de $x$, por lo 
que su omisión no afecta la expresión resultante de $C_p$. $\beta$ es el ángulo de inclinación (\emph{pitch}) de las palas.
\\ 

\textbf{Hipótesis 1:} Se asume $\beta = 0$ (sin control de \textit{pitch}), lo cual
es consistente con el alcance de la tesis, restringido a la operación en la región 2
(véase Sección \ref{sec:restricciones}), donde el ángulo de \textit{pitch} se mantiene fijo en un valor
cercano a cero para maximizar la extracción de potencia, en lugar de utilizarse
activamente como variable de control --- a diferencia de la región 3, donde $\beta$
se regula para limitar la potencia extraída ---.
\\

Bajo la Hipótesis 1, la superficie $C_p(\lambda,\beta)$ definida por
\eqref{eq:cp_analytic}--\eqref{eq:lambda_i} se reduce a una curva unidimensional
$C_p(\lambda) := C_p(\lambda,0)$, la cual es la que se utiliza en el resto del
análisis y las simulaciones de esta tesis. Utilizando estos parámetros se calculó
numéricamente mediante MATLAB el punto de máxima extracción de potencia,
obteniéndose $\lambda^\star = 7.96$ y $C_p^{\text{max}} = 0.411$. Estos parámetros y el
valor resultante de $\lambda^\star$ y $C_p^{\text{max}}$ coinciden con los reportados en
\cite{cisneros2020}, cuyo modelo matemático y parámetros de turbina se retoman en
esta tesis (véase Sección 2.3.1). La figura \ref{fig:cp_lambda} muestra la curva
característica $C_p$ vs $\lambda$ para estos parámetros.

% Utilizando estos parámetros y asumiendo $\beta = 0$ (sin control de \textit{pitch}), 
% se calculó numéricamente mediante MATLAB el punto de máxima extracción de potencia, 
% obteniéndose $\lambda^\star = 7.96$ y $C_p^{\text{max}} = 0.411$. Esta hipótesis es consistente 
% con el alcance de la tesis, restringido a la operación en la región 2 (véase Sección 2.2), 
% donde el ángulo de \textit{pitch} se mantiene fijo en un valor cercano a cero para maximizar 
% la extracción de potencia, en lugar de utilizarse activamente como variable de control 
% --- a diferencia de la región 3, donde $\beta$ se regula para limitar la potencia extraída---. 
% Bajo esta hipótesis, la superficie $C_p(\lambda,\beta)$ definida por \eqref{eq:cp_analytic}--\eqref{eq:lambda_i} 
% se reduce a una curva unidimensional $C_p(\lambda) := C_p(\lambda,0)$, la cual es la que se utiliza en 
% el resto del análisis y las simulaciones de esta tesis. Estos parámetros y el valor resultante de 
% $\lambda^\star$ y $C_p^{\text{max}}$ coinciden con los reportados en \cite{cisneros2020}, cuyo modelo 
% matemático y parámetros de turbina se retoman en esta tesis (véase Sección 2.3.1). La figura 
% \ref{fig:cp_lambda} muestra la curva característica $C_p$ vs $\lambda$ para estos parámetros.

\begin{figure}[ht]
\centering
\includegraphics[width=0.75\textwidth]{Cp_vs_lambda.pdf}
\caption{Curva de $\lambda$ vs $C_p$ calculada con los parámetros del sistema. 
El punto óptimo $\lambda^\star = 7.96$ corresponde al máximo coeficiente de potencia 
$C_p^{\text{max}} = 0.411$.}
\label{fig:cp_lambda}
\end{figure}

\subsection{Subsistema eléctrico}

La parte eléctrica del modelo está compuesta por el generador eléctrico (PMSG). La elección del generador 
eléctrico depende de factores como la potencia requerida para ser extraída del viento, la eficiencia 
requerida, la cantidad de mantenimiento requerida, la fiabilidad y la 
longevidad del generador en comparación con el resto de los componentes y la vida útil general de la 
turbina, entre otros. El desarrollo de esta tesis se limita al modelado y control del generador en sí; 
no se incluye el modelado de la electrónica de potencia del lado de la red ni de una carga específica, 
ya que estos elementos están fuera del alcance definido en la Sección \ref{sec:restricciones}.
\\

Dos de los tipos de generadores eléctricos más usados en WECS son los Generadores de Inducción Doblemente 
Alimentados (DFIG, por sus siglas en inglés: \emph{Doubly-Fed Induction Generator}) 
y los Generadores Síncronos de Imanes Permanentes (PMSG por sus siglas en inglés: \emph{Permament Magnent 
Synchronous Generator}). Por un lado, los generadores tipo DFIG, como 
su nombre lo indica, presentan una doble alimentación de voltaje hacia el rotor y el estator. Este último 
puede ser 
conectado directamente a la red y al rotor a través de un convertidor \cite{dahiya2019}. Este tipo de 
generadores es el más usado en turbinas eólicas de velocidad variable
\cite{dahiya2019,soltani2013}.
Los PMSG, por su parte, generan a una frecuencia eléctrica variable que depende directamente de la 
velocidad del rotor, por lo que requieren de un convertidor de potencia de capacidad plena (\emph{full-scale 
power converter}) para acoplar el estator a la red, procesando la totalidad de la potencia generada 
\cite{swibki2020}. Adicionalmente, los PMSG tienen la característica de permitir el accionamiento 
directo (\emph{direct drive}) del generador, lo cual hace posible la eliminación del \emph{gearbox} de 
la mecánica del sistema. Esto se logra a partir de incluir un alto número 
de polos en el diseño del generador, lo que permite la operación a bajas velocidades del rotor \cite{soltani2013, swibki2020}.
\\

Dentro de las consideraciones generales de ambos generadores al momento de decidir entre uno u otro 
se debe mencionar: i) la mayor eficiencia de los PMSG al operar efectivamente 
en un rango más amplio de condiciones climáticas (los DFIG operan con mayor eficiencia a velocidades 
altas del viento), ii) los costos de construcción asociados a los PMSG, 
en particular por los materiales utilizados en su fabricación; y iii) la necesidad de una caja de 
cambios en los DFIG, lo que incrementa su costo e introduce el problema del 
desgaste de estos componentes con el tiempo, así como la necesidad de mantenimiento a los componentes. 
Ambos tipos de generadores están diseñados  para operar a velocidad variables y permiten el control de 
la potencia reactiva \cite{dahiya2019, swibki2020,Lei2006,feijoo2000}. 
\\

Las ecuaciones dinámicas son similares en ambos 
tipos de generadores como se puede ver en \cite{Lei2006} para DFIG y en \cite{cisneros2020} para PMSG. 
Las diferencias se dan principalmente en que DFIG se incluyen también los voltajes del rotor además del estator \cite{Lei2006}.
El posterior desarrollo matemático en esta tesis asume la elección de un Generador Síncrono de Imanes 
Permanentes (PMSG).
\\

Los generadores son sistemas eléctricos trifásicos en los que cada señal de voltaje, intensidad de corriente 
o flujo magnético se encuentra desfasada $120^\circ$ con respecto a la anterior. En este sentido, a cada fase se le 
asigna una letra para denotar sus respectivas magnitudes. Este sistema se conoce como el \textbf{marco de referencia}
\emph{abc}. A pesar de que este marco de referencia emplea las variables físicas reales del sistema ($v_{abc}$, $i_{abc}$ y $\phi_{abc}$), 
introduce una mayor complejidad al momento de diseñar el control debido a la existencia de términos dependientes del 
tiempo y cuyo comportamiendo es senoidal. 
\\

Para simplificar el análisis matemático y el diseño del control se utiliza un cambio de variable que permite transformar 
las variables del sistema del marco de referencia \emph{abc} a un nuevo marco de referencia \emph{dq0} (en lo sucesivo nos referiremos 
únicamente a este marco de referencia como \emph{dq} o coordenadas \emph{dq}). En \cite{krause2013} este cambio de variable se maneja como
una transformación a un \textbf{marco de referencia arbitrario}, la cuál podría considerarse una versión genérica que agrupa a las trasformadas
específicas de Park \cite{park1929} para distintas aplicaciones. En particular, el cambio del marco de referencia \emph{abc} a \emph{dq} implica el uso de 
las transformadas de Clarke \cite{clarke1943} y Park de manera sucesiva. 
\\

En el nuevo marco de referencia las variables de tipo \emph{dq} se conocen como variables del \textbf{eje directo} para \emph{d} y variables 
del \textbf{eje de cuadratura} para \emph{q}. El eje directo se alinea con la dirección del flujo magnético producido por los imanes 
permanentes del rotor, mientras que el eje de cuadratura se encuentra en cuadratura (es decir, a 90 grados eléctricos) respecto al eje 
directo \cite{krause2013}. Esta elección de orientación es la que permite, como se detalla más adelante en este capítulo, desacoplar el 
control del flujo magnético del control del par electromagnético del generador. Las variables \emph{0's} no serán incluidas ya que dichas magnitudes están aritméticamente relacionadas
con las variables \emph{abc} y solo se incluirán aquellas que permitan condiciones de balance. i.e. \emph{dq} al ser independientes del tiempo \cite{krause2013,krause1998}.
\\

El análisis dinámico del generador de la turbina parte de leyes fundamentales de la electricidad
\footnote{El desarrollo completo de las ecuaciones dinámicas del sistema se puede encontrar en \cite{krause2013}}. 
En particular, para los \textbf{componentes resistivos} 
del generador, la ley de Ohm establece que:
\begin{align}
    v_{abc} = R_s i_{abc},
\end{align}
donde $v_{abc}, i_{abc} \in \mathbb{R}^3$ son los vectores de voltaje y corriente de fase, y $R_s \in \mathbb{R}^{3\times 3}$ 
es la matriz de resistencias del estator. Al transformar esta relación al marco de referencia arbitrario mediante la matriz 
de transformación $K_s$ --- que combina sucesivamente las transformadas de Clarke y Park, siendo esta última evaluada 
respecto al ángulo eléctrico del rotor $\theta_e = p\theta_m$, consistente con la orientación del eje directo hacia el 
flujo de los imanes permanentes discutida anteriormente --- obtenemos:
\begin{align}
    v_{dq0} = K_sR_s(K_s)^{-1} i_{dq0},
\end{align}
donde $v_{dq0}, i_{dq0} \in \mathbb{R}^3$ son los vectores transformados de voltaje y corriente.

Asumiendo que cada uno de los embobinados de fase del estator tienen la misma resistencia se cumple entonces que
\begin{align}
    K_sR_s(K_s)^{-1} = R,
\end{align}
donde $R \in \mathbb{R}$ es un escalar. Resultando en la siguiente expresión para el voltaje:
\begin{align}
    v_{dq} = R i_{dq}.
\end{align}

Por su parte, los \textbf{elementos inductivos} siguen la ley de inducción electromagnética de Faraday,
\begin{align}
    v_{abc} = \frac{d}{dt}\lambda_{abc},
\end{align}
donde $\lambda_{abc} \in \mathbb{R}^3$ son los enlaces de flujo, los cuales a su vez representan el 
flujo magnético total que atraviesa un embobinado y que, en el caso de un PMSG, recibe dos aportaciones principalmente:
el flujo magnético generado en los devanados del estator y el flujo producido en los imanes permanentes en el rotor.
Al transformar a coordenadas \emph{dq} obtenemos:
\begin{align}
    v_{dq} = \omega\lambda_{dq} + \frac{d}{dt}\lambda_{dq},
\end{align}
donde $\omega \in \mathbb{R}$ es la velocidad angular eléctrica del generador.

La expresión final de los voltajes \emph{dq} resulta de la suma de los voltajes de los elementos resistivos e inductivos. 
Por tanto,
\begin{align}
    v_{dq} = Ri_{dq} + \omega\lambda_{dq} + \frac{d}{dt}\lambda_{dq}.
\end{align}

En \cite{krause2013} los valores de $\lambda_{dq}$ se definen como sigue: 
\begin{align}
    &\lambda_q = L_q i_q,
    \\ 
    &\lambda_d = L_d i_d + \lambda_m,
\end{align}
donde $L_q = L_{ls} + L_{mq}$, $L_d = L_{ls} + L_{md}$ son escalares. i.e. en ambos casos, la inductancia 
total de cada eje es la suma de la inductancia de fuga más la inductancia de magnetización 
(inducida por los imanes permanentes); $\lambda_m \in \mathbb{R}$ es la constante 
del enlace de flujo magnético debido a los imanes permanentes.
\\

% El análisis dinámico del generador de la turbina parte de leyes fundamentales de la electricidad
% \footnote{El desarrollo completo de las ecuaciones dinámicas del sistema se puede encontrar en \cite{krause2013}}. 
% En particular, para los \textbf{componentes resistivos} 
% del generador, la ley de Ohm establece que:
% \begin{align}
%     v_{abc} = R_s i_{abc}
% \end{align}

% donde $R_s$ es la matriz de resistencias del estator. Al transformar esta relación al marco de referencia arbitrario obtenemos:
% \begin{align}
%     v_{dq0} = K_sR_s(K_s)^{-1} i_{dq0}
% \end{align}

% Asumiendo que cada uno de los embobinados de fase del estator tienen la misma resistencia se cumple entonces que:
% \begin{align}
%     K_sR_s(K_s)^{-1} = R
% \end{align}

% Resultando en la siguiente expresión para el voltaje:
% \begin{align}
%     v_{dq} = R i_{dq}
% \end{align}

% Por su parte, los \textbf{elementos inductivos} siguen la ley de inducción electromagnética de Faraday,
% \begin{align}
%     v_{abc} = \frac{d}{dt}\lambda_{abc}
% \end{align}

% donde $\lambda_{abc}$ son los enlaces de flujo, los cuales a su vez representan el 
% flujo magnético total que atraviesa un embobinado y que, en el caso de un PMSG, recibe dos aportaciones principalmente:
% el flujo magnético generado en los devanados del estator y el flujo producido en los imanes permanentes en el rotor.
% Al transformar a coordenadas \emph{dq} obtenemos:
% \begin{align}
%     v_{dq} = \omega\lambda_{dq} + \frac{d}{dt}\lambda_{dq}
% \end{align}

% La expresión final de los voltajes \emph{dq} resulta de la suma de los voltajes de los elementos resistivos e inductivos. 
% Por tanto,
% \begin{align}
%     v_{dq} = Ri_{dq} + \omega\lambda_{dq} + \frac{d}{dt}\lambda_{dq}
% \end{align}

% En \cite{krause2013} los valores de $\lambda_{dq}$ se asignan como sigue. 
% \begin{align}
%     &\lambda_q = L_q i_q
%     \\ 
%     &\lambda_d = L_d i_d + \lambda_m
% \end{align}

% donde $L_q = L_{ls} + L_{mq}$, $L_d = L_{ls} + L_{md}$. i.e. en ambos casos, la inductancia total de cada eje es la suma de la inductancia
% de fuga más la inductancia de magnetización (inducida por los imanes permanentes); $\lambda_m$ es el enlace de flujo de los imanes permanentes.
% \\

Una manipulación algebráica permite llegar a las siguientes ecuaciones dinámicas que caracterizan el funcionamiento 
del generador utilzando coordenadas \emph{dq}:


\begin{equation}
    \begin{aligned}
        L\frac{di_q}{dt} &= -Ri_d + pLi_q\omega_m - v_d
        \\
        L\frac{di_d}{dt} &= -Ri_q - pLi_d\omega_m + p\phi\omega_m - v_q 
    \end{aligned}
    \label{eq:electrical_subsystem}
\end{equation}

donde $i_d$, $i_q$, $v_d$ y $v_q$ son los voltajes y corrientes en el eje directo y de cuadratura, $p$ es el número de polos,
$\phi$ es el flujo magnético en el generador y $R$ es el valor de la resistencia en el generador. 
\\

\textbf{Hipótesis 2:} Para fines del diseño de control se asume que el sistema opera en la región 2.

\subsection{Sistema acoplado}

A partir de las ecuaciones del subsistema mecánico (\ref{eq:torque_balance_simplified}) y el subsistema eléctrico (\ref{eq:electrical_subsystem}) 
obtenemos el modelo dinámico que define el comportamiento del sistema completo de la turbina eólica:

\begin{equation}\label{eq:full_system}
    \begin{aligned}
        L\frac{di_q}{dt} &= -Ri_d + pLi_q\omega_m - v_d
        \\
        L\frac{di_d}{dt} &= -Ri_q - pLi_d\omega_m + p\phi\omega_m - v_q 
        \\
        J \frac{d\omega_m}{dt} &= -f_r \omega_m + T_m - T_e
    \end{aligned}
\end{equation}

Donde
\begin{align}
    T_m &:=T_m(\lambda)= r\frac{\kappa_1}{v_w} \frac{C_{p,0}(\lambda)}{\lambda}
    \\
    T_e  &:= T_e(i_q)= \frac{3}{2}p\phi i_q
    \\
    \lambda &:=\lambda(\omega_m)= \frac{r\omega_m}{v_w}
    \label{eq:definitions}
\end{align}

Donde $C_{p,0}(\lambda) := C_p(\lambda,0)$ es el coeficiente de potencia fijando el ángulo 
de inclinación en $\beta = 0$. Esta asignación de $\beta$ se debe a que el aerogenerador opera en la región 2 
y al objetivo de maximizar el coeficiente de potencia ($C_p$).
Los voltajes $v_d$ y $v_q$ se pueden manipular de forma libre, es decir, corresponden a las variables de control.

\section{Diseño del control de torque}

\subsection{Objetivos de control}
\label{sec:objetivos_de_control}

En turbinas éolicas de velocidad variable obtenemos una mayor eficiencia del sistema al reducir las pérdidas
del generador para una carga dada. Dentro del marco de referencia \emph{dq} la corriente $i_q$ se emplea para
inducir el torque en el generador. Esto se puede apreciar en la ecuación del torque eléctrico (\ref{eq:electrical_torque})
en la cual la corriente $i_d$ no tiene ninguna contribución. Por lo anterior, y para minimizar pérdidas, 
buscamos llevar la corriente del eje directo a un valor de $i_d = 0$ \cite{chinchilla2006}, ya que al regular esta 
corriente a cero se reducen las pérdidas en forma de calor en el generador, las cuales son de la forma $I^2R$, 
donde $I$ es la magnitud de la corriente del estator y $R$ es la resistencia del generador definida previamente.
Por su parte, el valor de la corriente $i_q$ deberá ser controlado con el objetivo de maximizar la extracción 
de potencia del viento. Lo anterior tiene como objetivo final la regulación de la velocidad rotacional del 
rotor de la turbina. 
\\

Previamente se introdujo el coeficiente de potencia. En dicha función se cumple que
para cada velocidad del viento existe una velocidad angular óptima que maximiza $C_p(\lambda,\beta)$ \cite{Leidhold2002,nak2013}.
Dicha velocidad angular óptima se conoce como \textbf{punto de máxima potencia} o MPP (por sus siglas en inglés: 
\emph{Maximum Power Point}) y al método o estrategia de control para encontrar este punto se le conoce como
\textbf{seguimiento del punto de máxima potencia} o MPPT (Por sus siglas en inglés: \emph{Maximum
Power Point Tracking}). Aunque existen varios métodos para el MPPT, en \cite{nak2013} se 
clasifican en tres principales categorías: control de la relación de velocidad punta (TSR), búsqueda 
\emph{Hill-climbing} (HCS), también conocida como \emph{Perturb \& Observe} (P\&O); y, por último, control de 
retroalimentación de la señal de potencia (PSF).
\\

Se considerará en el diseño del control del torque un algoritmo P\&O (o HCS) para el control de la velocidad óptima del
rotor. La estrategia principal en este tipo de control consiste en tomar muestras de la salida de potencia 
continuamente y compararlas con la salida inmediata anterior. Dependiendo de si la salida actual es mayor o menor
se incrementa o decremente el torque eléctrico del generador \cite{wang2004,buehring1981}.
\\

Una de las ventajas que ofrece un algoritmo HCS consiste en ser, de entre los tres tipos de estrategias
de control, el único que no requiere un conocimiento previo del sistema WECS y que es independiente de 
las características del viento, del generador y de la turbina en general. Lo anterior implica que el MPPT
se puede realizar en conjunto con otras estrategias de control para WECS como el \emph{pitch control} 
\cite{kazmi2011}. 
\\

Por otra parte, los algoritmos HCS presentan dos desventajas que pueden estar inducidas por alteraciones 
en la velocidad del viento en periodos cortos de tiempo o a un tamaño de paso del algoritmo inadecuado 
\cite{kazmi2011,xia2013}. Sin ambargo, estos problemas serán tratados con mayor detalle en la sección de
simulación dado que, como se mencionó, su implementación es independiente de la dinámica del generador y 
del resto de la turbina. 
\\

Por tanto, los objetivos de control se pueden resumir en:
\begin{enumerate}
    \item regular la corriente del eje directo en $i_d = i_d^\star = 0$ para minimizar pérdidas,
    \item regular la velocidad angular $\omega_m$ a un valor constante positivo que corresponda a 
        la velocidad óptima. Esta velocidad es impuesta por el algoritmo MPPT.
\end{enumerate}

\textbf{Hipótesis 3:} Para el diseño de los controladores se asume que el estado
completo del sistema --- la velocidad angular $\omega_m$ y las corrientes $i_d$,
$i_q$ --- se encuentra disponible para retroalimentación, ya sea de forma medida
o estimada. En la práctica, esto se logra mediante instrumentación estándar en
aerogeneradores con PMSG (por ejemplo, un encoder o tacómetro para $\omega_m$ y
sensores de corriente en los devanados del estator para $i_d$, $i_q$); el diseño
y la evaluación de dichos sensores quedan fuera del alcance de esta tesis, que se
enfoca en el análisis y diseño teórico del controlador.


\subsection{Estrategia de control}

El sistema de control propuesto combina distintas estrategias para regulación de la velocidad del rotor.
Para empezar, se regula el torque del generador mediante un control PI cuyo objetivo es mantener las variables
del PMSG en coordenadas \emph{dq} en un régimen adecuado para la máxima extracción de potencia del viento. Esta
estrategia corresponde al control por campo orientado (\emph{Field Oriented Control}, FOC), una de las técnicas
más empleadas para el control de torque en máquinas eléctricas, y cuya adopción para esta tesis fue justificada
en la Sección \ref{sec:soluciones_alternativas} frente a otras alternativas de la literatura
\cite{chinchilla2006,cisneros2020}. Para esto, se utiliza como referencia el valor de la velocidad rotacional y un
algoritmo de seguimiento del punto de máxima potencia (MPPT) para ajustar la velocidad rotacional óptima a la
cual se quiere seguir.
\\

Con base en la sección anterior, en particular las ecuaciones (\ref{eq:full_system}) que describen el 
sistema completo, nos interesa encontrar las ecuaciones que definen las variables de control $v_d$ y 
$v_q$. La estrategia utilizada para lograr esto es buscar un patrón de ``cancelación'' de términos 
no lineales combinado con la inserción de términos del control PI. De este modo, partiendo del modelo 
dinámico, se tiene para $v_d$:
\begin{align}
    v_d = pL i_q \omega_m + u_{\tt pi}^d \label{valUd},
\end{align}

donde $u^d_{\tt pi}$ es la salida del controlador PI definida, por tanto, como:
\begin{equation}
\begin{aligned}
    u_{\tt pi}^d &= k_{P1}(i_d-i_{d}^{\star}) + k_{I1}\xi_1\\
    &=k_{P1}i_d + k_{I1}\xi_1\\
    \dot{\xi_1} &= i_d-i_d^{\star}= i_d.
\end{aligned}
\label{eqUd}
\end{equation}

De manera similar para $v_q$:
\begin{align}
v_q = -pL i_d \omega_m + p\phi \omega_m + u_{\tt pi }^q,
\label{eqnVq}
\end{align}

donde $u_{\tt pi}^q$ corresponde al controlador PI definido por las siguientes ecuaciones:
\begin{equation}
\begin{aligned}
u^q_{\tt pi} &= k_{P2}(i_q - i_{q}^{ref}) + k_{I2}\xi_2 \\
\dot{\xi}_2 &= i_q - i_{q}^{ref},
\end{aligned}
\label{PI2}
\end{equation}

con la señal $i_q^{ref}$ definida como :
\begin{equation}\label{xi}
\begin{aligned}
i_{q}^{ref} &= k_{P3}(\omega_m - \omega_m^\star) + k_{I3}\xi_3 \\
\dot{\xi}_3 &= \omega_m - \omega_m^\star.
\end{aligned}
%\label{PI3}
\end{equation}

La estructura completa del control está definida por tres lazos de retroalimentación caracterizados
como sigue:

\begin{enumerate}
    \item \textbf{control de la corriente $i_d$}: \begin{itemize}
        \item entrada: la corriente $i_d$ medida que se busca mantener en 0,
        \item salida: el valor de $u_d$ para que $v_d$ compense las variaciones en $i_d$;
    \end{itemize}
    \item \textbf{control de la corriente $i_q$}: \begin{itemize}
        \item entrada: error de la corriente en el eje de cuadratura $(i_q - i_q^{ref})$,
        \item salida: el valor de $u_q$ para $v_q$;
    \end{itemize}
    \item \textbf{control de velocidad}: \begin{itemize}
        \item entrada: el error de la velocidad del rotor $(\omega_m - \omega_m^\star)$;
        \item salida: la señal de referencia de la corriente $i_q^{ref}$.
    \end{itemize}
\end{enumerate}

La sustitución de \eqref{eqUd} en \eqref{valUd} y, la sustitución de la expresión resultante en la primera
ecuación de \eqref{eq:electrical_subsystem} da lugar a las ecuaciones de $\Sigma_1$. Esto es:
\begin{align}
    L\frac{di_d}{dt} &= -(R+k_{P1})i_d - k_{I1}\xi_1\\
    \dot{\xi_1} &= i_d,
\end{align}

o, equivalentemente, % en representación de estados,
\begin{align}\label{eq:sigma_1}
    \mathbf{\Sigma_1:}\;\;\: \dot x= & \underbrace{\begin{bmatrix}0&1\\  -\frac{1}{L}k_{I1} & -\frac{1}{L}(R+k_{P1}) 
    \end{bmatrix}}_{A_1}x
 \end{align}

 con $x=\mathrm{col}(\xi_1,i_d)$. Claramente este subsistema corresponde a un sistema lineal e invariante en el tiempo. \\

Para el segundo subsistema, sustituyendo \eqref{eqnVq}, \eqref{PI2} y \eqref{xi} en la segunda ecuación de  \eqref{eq:electrical_subsystem}  resulta en la siguiente dinámica
\begin{align*}
    L\frac{di_q}{dt} =& -(R+k_{P2})i_q + k_{P2}i_q^{ref} \\
    =&-(R+k_{P2})i_q  + k_{P2}k_{P3}(\omega_m - \omega_m^\star) + k_{P2}k_{I3}\xi_3 - k_{I2}\xi_2.
\end{align*}

Uniendo la ecuación anterior con las de la dinámica de $\xi_2$ y de $\xi_3$ (descritas en  \eqref{PI2} y 
\eqref{xi}) resulta en la siguiente representación en espacio de estados
\begin{equation}\label{eq:sigma_2}
\resizebox{\textwidth}{!}{$
\mathbf{\Sigma_2:}\;\;\begin{aligned}
	\dot z=& \begin{bmatrix}
    \dot{\xi}_2 \\
    \dot{\xi}_3 \\
    \dot{\omega}_m \\
    \dot{i}_q
    \end{bmatrix} = \begin{bmatrix}
    z_4 - (k_{P3}(z_3 - \omega^\star_m) + k_{I3}z_2) \\
    z_3 - \omega^\star_m \\
    -\frac{f_r}{J}z_3 + \frac{r\kappa_1}{Jv_w}\frac{C_{p,0}(\lambda)}{\lambda} - \frac{3}{2J}p\phi z_4 \\
    -\frac{R+k_{P2}}{L}z_4 + \frac{k_{P2}k_{P3}}{L}(z_3 - \omega^\star_m) + \frac{k_{P2}k_{I3}}{L}z_2 - \frac{k_{I2}}{L}z_1
    \end{bmatrix}
\end{aligned}=:f(z)
$}
\end{equation}
donde $z=\mathrm{col}(\xi_2,\xi_3,\omega_m,i_q)$. Asumiendo la velocidad del 
viento constante,\footnote{Se aclara al lector que esta suposición se asume solamente para el propósito del control 
propuesto. Como se verá en las simulaciones, el hacer esto resulta en un control sencillo con un desempeño aceptable 
cuando el sistema es sometido a perfiles de viento realistas.} el sistema $\Sigma_2$ anterior es un sistema invariante del tiempo. 

% --------------- LEMA ------------------
La discusión anterior se resume en el siguiente lema.\\

\begin{lemma}\label{lem:wecs}
    Considere el sistema WECS \eqref{eq:full_system}-\eqref{eq:definitions} en lazo cerrado con \eqref{valUd}-\eqref{xi}. El sistema 
    resultante se compone de dos subsistemas desacoplados, que denominaremos $\Sigma_1$ y $\Sigma_2$, 
    cuyas ecuaciones son las siguientes:
    \begin{align*}
        \Sigma_1: \dot x= & \begin{bmatrix}0&1\\  -\frac{1}{L}k_{I1} & -\frac{1}{L}(R+k_{P1}) 
            \end{bmatrix}\\
        \Sigma_2: \dot z=& \begin{bmatrix}
            z_4 - (k_{P3}(z_3 - \omega^\star_m) + k_{I3}z_2) \\
            z_3 - \omega^\star_m \\
            -\frac{f_r}{J}z_3 + \frac{r\kappa_1}{Jv_w}\frac{C_{p,0}(\lambda)}{\lambda} - \frac{3}{2J}p\phi z_4 \\
            -\frac{R+k_{P2}}{L}z_4 + \frac{k_{P2}k_{P3}}{L}(z_3 - \omega^\star_m) + \frac{k_{P2}k_{I3}}{L}z_2 - \frac{k_{I2}}{L}z_1
            \end{bmatrix}
    \end{align*}
    \tripletriang
\end{lemma}

A continuación, se linealizará el sistema alrededor del punto de equilibrio que
cumpla con los objetivos de control. Luego, se obtendrán las condiciones en las
ganancias $k_P$s y $k_I$s que hacen que esa linealización tenga valores propios en
el semiplano izquierdo del plano complejo. De acuerdo con el Método Indirecto de
Lyapunov \cite{khalil2002}, esto garantiza que el punto de equilibrio
deseado sea exponencialmente estable.


\subsubsection*{Puntos de equilibrio de $\Sigma_1$ y $\Sigma_2$}

Como parte del análisis de estabilidad, primero se obtienen los puntos de equilibrio de ambos subsistemas. 
Estos quedan enunciados en los siguientes lemas.

\begin{lemma}
    El sistema $\Sigma_1$ tiene un solo punto de equilibrio en $$x=\bar{x}:=0$$
    \tripletriang

    \begin{proof}
        Al igualar $\dot x=0$ en (1.37), resulta en la siguientes ecuaciones
        \begin{align*}
            \bar{x}_2=&0\\
            -(R+k_{P1})x_2-k_{I1}x_1  =&0 	
        \end{align*}
        Se concluye que $\bar{x}_2=0$ es un punto de equibrio. Reemplazando lo anterior en la segunda ecuación queda
        $$k_{I1}x_1  =0,$$
        de donde se concluye que  tambien  $x_1=0$ es el punto de equilibrio para $x_1$.
    \end{proof}
\end{lemma}

\begin{lemma}\label{lem:sigma_2_eq}
    El sistema $\Sigma_2$ tiene un único punto de equilibrio dado por
	\begin{align}\label{barz}
	\bar{z} :=
		\begin{bmatrix}
			-\frac{2R}{3p\phi k_{\textrm{I2}}}(T^\star_m - f_r\omega^\star_m) \\[0.3cm]
			\frac{2}{3p\phi k_{\textrm{I3}}}(T^\star_m - f_r\omega^\star_m) \\[0.3cm]
			\omega^\star_m \\[0.3cm]
			\frac{2}{3p\phi}(T^\star_m - f_r\omega^\star_m)
	\end{bmatrix}
	\end{align}
    donde $T^\star_m := r\frac{\kappa_1 v_w C_{p,0}(\lambda^\star)}{\lambda^\star}$ es el torque mecánico en el 
    equilibrio, $\lambda^\star := \frac{r\omega_m^\star}{v_w}$ y $\omega_m^\star$ es la referencia de velocidad 
    angular.

    \tripletriang

    \begin{proof}
        Los puntos de equilibrio se obtienen al igualar a cero las ecuaciones del sistema \eqref{eq:sigma_2}, esto es, resolviendo $\dot{z} = 0$:

        \begin{align}
        z_4 - [k_{\textrm{P3}}(z_3 - \omega^\star_m) + k_{\textrm{I3}}z_2] &= 0 \label{eq:ss1} \\
        z_3 - \omega^\star_m &= 0 \label{eq:ss2} \\
        -\frac{f_r}{J}z_3 + \frac{r\kappa_1}{Jv_w}\frac{C_{p,0}(\lambda)}{\lambda} - \frac{3}{2J}p\phi z_4 &= 0 \label{eq:ss3} \\
        -\frac{R+k_{\textrm{P2}}}{L}z_4 + \frac{k_{\textrm{P2}}k_{\textrm{P3}}}{L}(z_3 - \omega^\star_m) + \frac{k_{\textrm{P2}}k_{\textrm{I3}}}{L}z_2 - \frac{k_{\textrm{I2}}}{L}z_1 &= 0 \label{eq:ss4}
        \end{align}

        De la ecuación \eqref{eq:ss2} se obtiene directamente que:
        \begin{equation}
        \bar{z}_3 = \omega^\star_m \label{eq:eq1}
        \end{equation}

        Sustituyendo \eqref{eq:eq1} en \eqref{eq:ss1}, se tiene:
        \begin{align}
        \bar{z}_4 - k_{\textrm{I3}}\bar{z}_2 = 0 \label{eq:eq2}
        \end{align}

        Por otro lado, de la ecuación \eqref{eq:ss3}, usando \eqref{eq:eq1} y $\lambda^\star = \frac{r\omega^\star_m}{v_w}$, se obtiene:
        \begin{equation}
        \bar{z}_4 = \frac{2}{3p\phi}(T^\star_m - f_r\omega^\star_m) \label{eq:eq3}
        \end{equation}

        Combinando \eqref{eq:eq2} y \eqref{eq:eq3}:
        \begin{equation}
        \bar{z}_2 = \frac{1}{k_{\textrm{I3}}}\bar{z}_4 = \frac{2}{3p\phi k_{\textrm{I3}}}(T^\star_m - f_r\omega^\star_m) \label{eq:eq4}
        \end{equation}

        Finalmente, sustituyendo $\bar{z}_3 = \omega^\star_m$ y \eqref{eq:eq3} en \eqref{eq:ss4}:
        \begin{equation}
        \bar{z}_1 = -\frac{R}{k_{\textrm{I2}}}\bar{z}_4 = -\frac{2R}{3p\phi k_{\textrm{I2}}}(T^\star_m - f_r\omega^\star_m) \label{eq:eq5}
        \end{equation}

        Así, el punto de equilibrio queda determinado por \eqref{barz}.
    \end{proof}
\end{lemma}

\subsection{Análisis de estabilidad de \texorpdfstring{$\Sigma_1$ y $\Sigma_2$}{Sigma 1 y Sigma 1}}
\label{sec:analisis_sigmas}

Una vez que se encuentran los puntos de equilibrio, en esta subsección se analizará la estabilidad de 
los mismos. Puesto que los subsistemas están desacoplados, la estabilidad de cada uno se puede estudiar 
por separado.\\

\begin{lemma}\label{lem:sigma_1}
    Considere el sistema $\Sigma_1$ con $k_{P1}>-R$ y $k_{I1}>0$. El punto de equilibrio 
    $x=\bar x=0$ es exponencialmente, localmente estable.

    \tripletriang

    \begin{proof}
        El polinomio característico de $A_1$ en (\ref{eq:sigma_1}) es el siguiente:
        \begin{align}\label{p1}
            p_1(s)=s^2 + \frac{k_{P1}+R}{L}s + \frac{k_{I1}}{L}.
        \end{align}
    
        Para un polinomio de grado dos, la tabla de Routh--Hurwitz asociada a
        $p_1(s) = a_2 s^2 + a_1 s + a_0$ tiene como primera columna a los
        coeficientes $a_2$, $a_1$ y $a_0$, por lo que el criterio de estabilidad de
        Routh se reduce, en este caso particular, a la condición equivalente de que
        todos los coeficientes del polinomio sean positivos. Por tanto, tomando en
        cuenta que $L>0$, tenemos las siguientes relaciones:
        \begin{align}
            k_{P1} +R > 0  \Longleftrightarrow k_{P1} &> -R,\\
            k_{I1} &> 0.
        \end{align}
    \end{proof}
\end{lemma}

% -------------------\\
% De lo anterior, los valores proporcional e integral $k_{P1} = 0$ y $k_{I1} = -1$ garantizan la estabilidad de $\Sigma_1$.\\
% -------------------

Por otro lado, para estudiar la estabilidad del subsistema $\Sigma_2$, se procederá a linealizarlo.
Para tal motivo se calcula el Jacobiano y se evalúa en  el punto de equilibrio. Esto es, el sistema 
linealizado tiene la forma:
\begin{align}\label{dz}
\Sigma_{2}^\mathrm{lin}:\;\;\;\;\dot{\delta z} =A_2 \delta z
\end{align}
donde $\delta z$ representa al estado del sistema linealealizado,  
\begin{equation}
A_2 = \left.\frac{\partial f}{\partial z}\right|_{z=\bar{z}} = \begin{bmatrix}
0 & -k_{I3} & -k_{P3} & 1 \\
0 & 0 & 1 & 0 \\
0 & 0 & -\frac{f_r}{J} + \gamma & -\frac{3}{2J}p\phi \\
-\frac{k_{I2}}{L} & \frac{k_{P2}k_{I3}}{L} & \frac{k_{P2}k_{P3}}{L} & -\frac{R+k_{P2}}{L}
\end{bmatrix} \label{eq:A2_symbolic}
\end{equation}

donde $\gamma$ está definida como la siguiente derivada evaluada en el punto de
equilibrio $\lambda=\lambda^\star$ (véase la Sección \ref{sec:coef_potencia}, donde se define
$\lambda^\star$ como el valor de la relación de velocidad punta óptima)
\begin{align}\nonumber
	\gamma:= \frac{r\kappa_1}{Jv_w}\left.\frac{\partial}{\partial \lambda}
	\left(\frac{C_{p,0}(\lambda)}{\lambda}\right)\right|_{\lambda=\lambda^\star}.
\end{align}
Desarrollando la derivada, la constante $\gamma$ se reduce a la expresión
\begin{align}
	\gamma=&\left.\frac{r\kappa_1}{Jv_w}
	\left[\frac{1}{\lambda}\frac{\partial }{\partial \lambda}C_{p,0}(\lambda)
	-\frac{1}{\lambda^2}C_{p,0}(\lambda)\right]\right|_{\lambda=\lambda^\star}\nonumber\\
	=& - \frac{1}{(\lambda^\star)^2}\frac{r\kappa_1}{Jv_w} C_{p,0}(\lambda^\star). \label{eq:gamma}
\end{align}


Puesto que los parámetros son positivos y $C_{p,0}>0$, entonces $\gamma<0$. A partir de lo anterior, 
es posible enunciar el siguiente lema. Así mismo, es posible concluir que la entrada (3,3) de la 
matriz $A_2$ toma valores negativos, esto es, $A_{2}^{(3,3)}<0$.\\

\begin{lemma}\label{lem:sigma_2}
    Considere el sistema $\Sigma_2^{\mathrm{lin}}$ en \eqref{dz} que consiste en la
    linealización de $\Sigma_2$ alrededor de su punto de equilibrio. El
    punto de equilibrio $z=\bar{z}$ de $\Sigma_2$, con la constante $\bar{z}$
    definida en \eqref{barz}, es exponencial y localmente estable si y
    solo si se cumplen las siguientes condiciones:
    \begin{align}
        & \frac{R}{L} + \bar{\gamma_1} + k_{\textrm{P2}}' > 0 \label{eq:cond1}
        \\
        & \bar{\gamma_2}k_{\textrm{I2}}'k_{\textrm{I3}} > 0 \label{eq:cond2}
        \\
        & \varrho > 0 \label{eq:cond3}
        \\
        & \bar{\gamma_2}\left[k_{\textrm{I2}}'k_{\textrm{P3}} + k_{\textrm{I3}}k_{\textrm{P2}}'\right] + \bar{\gamma_1}k_{\textrm{I2}}'
    - \frac{\bar{\gamma_2}k_{\textrm{I2}}'k_{\textrm{I3}}\left(\frac{R}{L} + \bar{\gamma_1} + k_{\textrm{P2}}'\right)}{\varrho} > 0 \label{eq:cond4}
    \end{align}
    
    donde 
    \begin{align*}
        \varrho := k_{\textrm{P2}}' \left[\bar{\gamma_1} + \bar{\gamma_2}k_{\textrm{P3}}\right] + \frac{R}{L}\bar{\gamma_1} + k_{\textrm{I2}}' 
    - \frac{k_{\textrm{I2}}'\left[\bar{\gamma_1} + \bar{\gamma_2}k_{\textrm{P3}}\right] + \bar{\gamma_2}k_{\textrm{I3}}k_{\textrm{P2}}'}{\frac{R}{L} + \bar{\gamma_1} + k_{\textrm{P2}}'}
    \end{align*}
    
    con $k_{\textrm{P2}}' := \frac{k_{\textrm{P2}}}{L}$, $k_{\textrm{I2}}' := \frac{k_{\textrm{I2}}}{L}$, $\bar{\gamma}_1 := \frac{f_r}{J} - \frac{\gamma}{J} > 0$ y $\bar{\gamma}_2 := \frac{3}{2J}p\phi > 0$.

    \tripletriang
    \begin{proof}
        Bajo las definiciones de $k_{\textrm{I2}}'$, $k_{\textrm{P2}}'$, $\bar{\gamma}_1$ y $\bar{\gamma}_2$ que se muestran en el enunciado del lema, la matriz $A_2$ de $\Sigma_2^{\mathrm{lin}}$ es
        $$\bar{A}_2 := \begin{bmatrix}
            0 & -k_{\textrm{I3}} & -k_{\textrm{P3}} & 1 \\
            0 & 0 & 1 & 0 \\
            0 & 0 & -\bar{\gamma}_1  & -\bar{\gamma}_2 \\
            -k_{\textrm{I2}}' & k_{\textrm{P2}}'k_{\textrm{I3}} & k_{\textrm{P2}}'k_{\textrm{P3}} & -\left(\frac{R}{L}+k_{\textrm{P2}}'\right)
        \end{bmatrix}$$

        El polinomio característico de la matriz $\bar{A}_2$, calculado mediante MATLAB, es:

        \begin{equation}\label{eq:poly_sigma_2}
            \begin{split}
                p_2(s) =& s^4 + \left(\frac{R}{L} + \bar{\gamma_1} + k_{\textrm{P2}}'\right)s^3 + \ldots 
                \\ 
                &\ldots + \left( \frac{R\bar{\gamma_1}}{L} + k_{\textrm{I2}}' + \bar{\gamma_1}k_{\textrm{P2}}' + \bar{\gamma_2}k_{\textrm{P2}}'k_{\textrm{P3}} \right)s^2 + \ldots 
                \\
                &\ldots + \left( \bar{\gamma_1}k_{\textrm{I2}}' + \bar{\gamma_2}k_{\textrm{I2}}'k_{\textrm{P3}} + \bar{\gamma_2}k_{\textrm{I3}}k_{\textrm{P2}}' \right)s + \ldots
                \\
                &\ldots + \bar{\gamma_2}k_{\textrm{I2}}'k_{\textrm{I3}}
            \end{split}
        \end{equation}

        que tiene la forma general:
        \begin{equation}\label{eq:poly_sigma_2_form}
            p_2(s) = a_4s^4 + a_3s^3 + a_2s^2 + a_1s + a_0
        \end{equation}

        donde $a_4 = 1$ y los demás coeficientes se identifican por comparación con \eqref{eq:poly_sigma_2}.

        Por el criterio de Routh-Hurwitz, el sistema es estable si y solo si todos los elementos de la primera columna de la tabla de Routh son positivos. La tabla de Routh-Hurwitz tiene la siguiente estructura:

        \begin{equation}
        \begin{array}{c|ccc}
        s^4 & a_4 & a_2 & a_0 \\
        s^3 & a_3 & a_1 & \textcolor{white}{0} \\
        s^2 & b_1 & b_2 & \textcolor{white}{0} \\
        s^1 & c_1 & \textcolor{white}{0} & \textcolor{white}{0} \\
        s^0 & d_1 & \textcolor{white}{0} & \textcolor{white}{0}
        \end{array} \label{eq:rh_table}
        \end{equation}

        donde los elementos de la primera columna se calculan como:
        \begin{align}
        b_1 &= \frac{a_3a_2 - a_4a_1}{a_3} \label{eq:b1} \\
        c_1 &= \frac{b_1a_1 - a_3b_2}{b_1} \label{eq:c1} \\
        d_1 &= b_2 = a_0 \label{eq:d1}
        \end{align}

        Las condiciones de estabilidad requieren que:
        \begin{equation}\label{eq:routh_conditions}
            a_4 > 0, \quad a_3 > 0, \quad b_1 > 0, \quad c_1 > 0, \quad d_1 > 0
        \end{equation}

        Dado que $a_4 = 1$, esta condición se satisface automáticamente.

        La condición $a_3 > 0$ conduce directamente a \eqref{eq:cond1}:
        \begin{equation*}
            \frac{R}{L} + \bar{\gamma_1} + k_{\textrm{P2}}' > 0
        \end{equation*}

        La condición $d_1 = a_0 > 0$ resulta en \eqref{eq:cond2}. Dado que $\bar{\gamma}_2 > 0$ por definición, se requiere:
        \begin{equation*}
            k_{\textrm{I2}}'k_{\textrm{I3}} > 0
        \end{equation*}

        Para la condición $b_1 > 0$, sustituyendo $a_4 = 1$ en \eqref{eq:b1}:
        \begin{equation*}
            b_1 = \frac{a_3a_2 - a_1}{a_3} = a_2 - \frac{a_1}{a_3} > 0
        \end{equation*}

        Desarrollando esta expresión con los valores explícitos de $a_1$, $a_2$ y $a_3$ se obtiene la condición \eqref{eq:cond3}: $\varrho > 0$.

        Finalmente, para la condición $c_1 > 0$, notando que $b_2 = a_0$:
        \begin{equation*}
            c_1 = \frac{b_1a_1 - a_3a_0}{b_1} = a_1 - \frac{a_3a_0}{b_1} > 0
        \end{equation*}

        Sustituyendo los valores de $a_0$, $a_1$ y $a_3$, y usando la expresión de $b_1$ en términos de $\varrho$, se obtiene la condición \eqref{eq:cond4}.
    \end{proof}
\end{lemma}

\subsection*{Algoritmo para encontrar las ganancias}\label{sec:algoritmo}

El siguiente algoritmo proporciona un procedimiento sistemático para la 
sintonización de un conjunto de ganancias que satisfagan las condiciones 
\eqref{eq:cond1}--\eqref{eq:cond4}.

\begin{tabular}{l}
\hline
\textbf{Algoritmo.} Encontrar ganancias $k_{\textrm{P2}}',k_{\textrm{P3}},k_{\textrm{I2}}',k_{\textrm{I3}}$\\
\hline
\textbf{Paso 1}. Asignar valores positivos a $k_{\textrm{P3}},k_{\textrm{I3}}$.\\
\textbf{Paso 2}. Fijar $k_{\textrm{I2}}' = 1/k_{\textrm{P2}}'$\\
\textbf{Paso 3.} Encontrar $k_{\textrm{P2}}' > 0$ suficientemente grande tal que se satisfagan \\
\quad\quad\quad\quad las condiciones \eqref{eq:cond3} y \eqref{eq:cond4}\\\hline
\end{tabular}
\vspace{1em}

Obsérvese que, de acuerdo con el Paso 1 y Paso 2, al hacer $k_{\textrm{P3}} = \bar{k}_{\textrm{P3}} > 0$, $k_{\textrm{I3}} = \bar{k}_{\textrm{I3}} > 0$ y $k_{\textrm{I2}}' = \frac{1}{k_{\textrm{P2}}'}$, la condición \eqref{eq:cond1} se satisface automáticamente para $k_{\textrm{P2}}' > 0$ suficientemente grande. Asimismo, la condición \eqref{eq:cond2} se cumple por la elección de ganancias positivas. Por lo tanto, solo resta verificar las condiciones \eqref{eq:cond3} y \eqref{eq:cond4}.

Sustituyendo estos valores en \eqref{eq:cond3}, la desigualdad se reduce a:
\begin{align*}
&\varrho\left(k_{\textrm{P2}}',\bar{k}_{\textrm{P3}},\frac{1}{k_{\textrm{P2}}'},\bar{k}_{\textrm{I3}}\right) = \\
&\quad k_{\textrm{P2}}'\left[\bar{\gamma}_1 + \bar{\gamma}_2 \bar{k}_{\textrm{P3}}\right] + \frac{R}{L}\bar{\gamma}_1  + \frac{1}{k_{\textrm{P2}}'} - \frac{\frac{1}{k_{\textrm{P2}}'^2}\left[\bar{\gamma}_2\bar{k}_{\textrm{P3}} + \bar{\gamma}_1\right] + \bar{\gamma}_2 \bar{k}_{\textrm{I3}}}{1 + \frac{1}{L}\frac{R}{k_{\textrm{P2}}'} + \frac{\bar{\gamma}_1}{k_{\textrm{P2}}'}} > 0
\end{align*}

A medida que $k_{\textrm{P2}}'$ toma valores más grandes, el término dominante $k_{\textrm{P2}}'\left[\bar{\gamma}_1 + \bar{\gamma}_2 \bar{k}_{\textrm{P3}}\right]$ garantiza que \eqref{eq:cond3} se satisface. Además, un valor creciente de $k_{\textrm{P2}}'$ implica un valor creciente de $\varrho$, lo cual facilita el cumplimiento de la condición \eqref{eq:cond4}:
\begin{align}
\bar{\gamma}_2 \left[\frac{1}{k_{\textrm{P2}}'}\bar{k}_{\textrm{P3}} + \bar{k}_{\textrm{I3}}k_{\textrm{P2}}'\right] + \frac{\bar{\gamma}_1}{k_{\textrm{P2}}'} - \frac{\frac{1}{k_{\textrm{P2}}'}\left(\frac{R}{L} + k_{\textrm{P2}}' + \bar{\gamma}_1\right)\bar{\gamma}_2\bar{k}_{\textrm{I3}}}{\varrho\left(k_{\textrm{P2}}',\bar{k}_{\textrm{P3}},\frac{1}{k_{\textrm{P2}}'},\bar{k}_{\textrm{I3}}\right)} > 0
\end{align}

ya que el término $\bar{k}_{\textrm{I3}}k_{\textrm{P2}}'$ crece linealmente mientras 
que el término del denominador decrece. Por lo tanto, existe un valor de 
$k_{\textrm{P2}}'$ suficientemente grande que garantiza que las condiciones 
de Routh-Hurwitz se satisfagan. Del análisis anterior se puede hacer la siguiente
aserción.

\begin{lemma}\label{lem:conds}
    Las condiciones \eqref{eq:cond1}--\eqref{eq:cond4} admiten solución. \tripletriang
\end{lemma}

